\subsection{مبانی نظری و فرمول‌بندی ریاضی}
فرض کنید $\mathbf{W} \in \mathbb{R}^{d_{\text{out}} \times d_{\text{in}}}$ ماتریس وزن یک لایه خطی متراکم و $\mathbf{X} \in \mathbb{R}^{B \cdot N \times d_{\text{in}}}$ تانسور فعال‌سازی‌های ورودی استخراج‌شده از داده‌های کالیبراسیون باشد. بردار انرژی فعال‌سازی با گشتاور دوم به صورت زیر محاسبه می‌شود:
\begin{equation}
\mathbf{s}_j = \sqrt{\mathbb{E}[x_j^2] + \epsilon} = \sqrt{\frac{1}{M}\sum_{i=1}^M \mathbf{X}_{i, j}^2 + \epsilon}, \quad \forall j \in \{1, \dots, d_{\text{in}}\}
\end{equation}
سپس ماتریس وزن‌ها با قطر این بردار مقیاس‌گذاری می‌شود تا ویژگی‌های پرانرژی تقویت گردند:
\begin{equation}
\tilde{\mathbf{W}} = \mathbf{W} \operatorname{diag}(\mathbf{s})
\end{equation}
بر روی ماتریس مقیاس‌یافته، تجزیه مقادیر منفرد فشرده اجرا می‌شود:
\begin{equation}
\tilde{\mathbf{W}} = \mathbf{U} \mathbf{\Sigma} \mathbf{V}^\top
\end{equation}
فاکتورهای کم‌مرتبه متناظر با رتبه انتخابی $r \ll \min(d_{\text{in}}, d_{\text{out}})$ با توزیع متقارن مقادیر منفرد و اعمال معکوس مقیاس ورودی تشکیل می‌شوند:
\begin{equation}
\mathbf{U}_r = \mathbf{U}_{:, :r} \mathbf{\Sigma}_{:r}^{1/2}, \quad \mathbf{V}_r = \mathbf{\Sigma}_{:r}^{1/2} \mathbf{V}_{:, :r}^\top \operatorname{diag}(\mathbf{s})^{-1}
\end{equation}
در نتیجه بازسازی موثر وزن‌ها به فرم فاکتورگیری‌شده کم‌مرتبه درمی‌آید:
\begin{equation}
\hat{\mathbf{W}} = \mathbf{U}_r \mathbf{V}_r \approx \mathbf{W}
\end{equation}
در گام پایانی، کالیبراسیون مجدد لایه‌های نرمال‌سازی به منظور هماهنگ‌سازی واریانس فعال‌سازی‌های جدید صورت می‌گیرد:
\begin{equation}
\boldsymbol{\gamma}_{\text{new}} = \boldsymbol{\gamma}_{\text{old}} \odot \frac{\boldsymbol{\sigma}_{\text{new}}}{\boldsymbol{\sigma}_{\text{old}}}
\end{equation}

که در آن $\sigma_i(\mathbf{W})$ بیانگر $i$-امین مقدار منفرد مرتب‌شده ماتریس به صورت نزولی است. نسبت حفظ انرژی طیفی $\eta_r = \frac{\sum_{i=1}^r \sigma_i^2}{\sum_{i=1}^{d} \sigma_i^2}$ تضمین می‌کند که با انتخاب آستانه $\tau \ge 0.90$، بیش از نود درصد از پراکندگی واریانس داده‌ها حفظ شده و افت سیگنال در انتشار گرادیان رو به عقب به حداقل ممکن کاهش می‌یابد.

\paragraph{تحلیل مجانبی و پیچیدگی اشتراک حافظه سخت‌افزاری}
از منظر پیچیدگی محاسباتی و انتقال داده‌ها در سلسومراتب حافظه (حافظه سراسری \lr{HBM} به حافظه محلی \lr{SRAM})، عملیات ضرب ماتریسی در این مکانیزم نیازمند تخصیص بافر موقت با ابعاد فاکتورگیری‌شده است. برای پردازش دنباله‌ای به طول $N$ با بعد مخفی $D$، بار محاسباتی هسته از فرم درجه دو به رابطه زیر تقلیل می‌یابد:
\begin{equation}
\label{eq:complexity_a3_compressor}
\mathcal{C}_{\text{a3\_compressor}} = \mathcal{O}\left( B \cdot N \cdot D \cdot r + B \cdot H \cdot N \cdot d_k \right) \ll \mathcal{O}\left( B \cdot H \cdot N^2 \cdot d_k \right)
\end{equation}
این بهبود مرتبه پیچیدگی، با کاهش چشمگیر حجم تبادل تانسورها در گذرگاه‌های ارتباطی کارت گرافیک، مانع از اشباع پهنای باند حافظه (پدیده \lr{Memory-Bound}) در فرآیند استنتاج طولانی‌مدت می‌گردد.
